\subsection{Dynamic Pricing Mechanism}
\label{sec:pricing}

BPE's backpressure allocation determines \emph{who} receives payment; we now
formalize the complementary question of \emph{how much} each unit of service
costs. We introduce a dynamic pricing curve that creates \emph{economic
backpressure}: as queue load builds at a sink, the per-unit price rises,
naturally throttling demand before overflow buffers are needed.

\begin{definition}[Queue-Length Price]
\label{def:price}
For task type $\tasktype$ and sink $k$ at time $t$, the per-unit price is:
\begin{equation}
\label{eq:price}
p(\tasktype, k, t) = \beta_\tasktype(t) \cdot
    \left(1 + \gamma \cdot \frac{Q(\tasktype, k, t)}{\Csmooth(k, t, \tasktype)}\right)
\end{equation}
where $\beta_\tasktype(t) > 0$ is the base fee for task type $\tasktype$,
$\gamma > 0$ is the price sensitivity parameter,
$Q(\tasktype, k, t) \geq 0$ is the current queue load reported by sink $k$,
and $\Csmooth(k, t, \tasktype)$ is the EWMA-smoothed capacity from \Cref{sec:ewma}.
\end{definition}

When a sink has zero queue load, the price reduces to the base fee $\beta_\tasktype$.
As $Q / \Csmooth \to 1$ (queue approaches capacity), the price doubles.
When $Q > \Csmooth$ (overloaded), the price rises further, discouraging
additional demand.

\paragraph{Base fee adjustment.}
The base fee $\beta_\tasktype$ adjusts per epoch following an EIP-1559-style
mechanism~\cite{roughgarden2021eip1559}. Let $D(\tasktype, t)$ denote aggregate
demand (total flow rate from sources) and $C_{\text{total}}(\tasktype, t) =
\sum_k \Csmooth(k, t, \tasktype)$ be aggregate capacity. At each epoch boundary:

\begin{equation}
\label{eq:basefee}
\beta_\tasktype(t+1) =
\begin{cases}
\beta_\tasktype(t) \cdot (1 + \delta) & \text{if } D(\tasktype, t) > C_{\text{total}}(\tasktype, t), \\
\beta_\tasktype(t) \cdot (1 - \delta) & \text{otherwise},
\end{cases}
\end{equation}
where $\delta = 0.125$ (12.5\%), matching the maximum per-block adjustment rate
of EIP-1559. A floor $\beta_{\min}$ prevents the base fee from collapsing to zero
during sustained under-utilization.

\begin{proposition}[Price Equilibrium]
\label{prop:price-eq}
Under the adjustment rule \eqref{eq:basefee} with price-sensitive demand
$D(\tasktype, \beta) = D_0 / \beta$ (unit-elastic), there exists a unique
equilibrium base fee:
\begin{equation}
\beta^* = \frac{D_0}{C_{\text{total}}(\tasktype)}
\end{equation}
at which $D(\tasktype, \beta^*) = C_{\text{total}}(\tasktype)$.
\end{proposition}

\begin{proof}
At equilibrium, the adjustment rule \eqref{eq:basefee} is stationary, requiring
$D(\tasktype, \beta^*) = C_{\text{total}}$. Substituting the demand function:
$D_0 / \beta^* = C_{\text{total}}$, yielding
$\beta^* = D_0 / C_{\text{total}}$. Uniqueness follows from
strict monotonicity of $D(\cdot)$ in $\beta$. Stability follows because
$\beta < \beta^*$ implies excess demand, triggering $\beta \uparrow$,
and $\beta > \beta^*$ implies excess capacity, triggering $\beta \downarrow$.
\end{proof}

\paragraph{Connection to Kelly shadow prices.}
\Cref{prop:price-eq} recovers a classical result from network pricing theory.
Kelly~\cite{kelly1998rate} showed that the Lagrange multiplier on each link's
capacity constraint in the proportional fairness optimization emerges as a
per-unit price. In BPE, the equilibrium base fee $\beta^*$ is precisely this
shadow price for the aggregate capacity constraint. The queue-length surcharge
$\gamma \cdot Q/\Csmooth$ in \eqref{eq:price} refines this by providing
per-sink price differentiation: sinks with shorter queues offer lower prices,
attracting more demand, which is exactly the routing signal that a
price-taking source should follow.

\begin{proposition}[Optimal Source Routing]
\label{prop:routing}
A cost-minimizing source facing prices \eqref{eq:price} across sinks routes
flow to the sink with the lowest price. At equilibrium, source flow distributes
such that prices equalize across sinks with available capacity:
\begin{equation}
Q(\tasktype, k, t) / \Csmooth(k, t, \tasktype) = Q(\tasktype, j, t) / \Csmooth(j, t, \tasktype)
\quad \forall\, k, j \in K_\tasktype
\end{equation}
i.e., all active sinks operate at the same utilization ratio.
\end{proposition}

\begin{proof}
If sink $k$ has a lower utilization ratio $Q_k/\Csmooth_k < Q_j/\Csmooth_j$,
then by \eqref{eq:price}, $p(\tasktype, k, t) < p(\tasktype, j, t)$.
A cost-minimizing source strictly prefers $k$, shifting demand until
utilization ratios equalize. At this point, no unilateral deviation reduces
cost, establishing a Nash equilibrium in routing.
\end{proof}

\paragraph{Economic backpressure.}
Traditional backpressure routing uses queue differentials as a signal to route
traffic away from congested links~\cite{tassiulas1992stability}. Our pricing
mechanism achieves the same effect through economic incentives: congested sinks
become expensive, and rational sources route elsewhere. The key advantage is
that pricing operates \emph{without centralized coordination}---each source
independently queries per-sink prices and routes to minimize cost, yet the
aggregate effect is optimal load balancing.
